\section*{Bab \ref{ch:b3:groups} --- Teori Grup}

\begin{solution}{exo:b3:groups:1}
(a) Misalkan $[G:H] = 2$. Untuk $g \in H$, $gH = H = Hg$. Untuk $g \notin
H$: kedua koset kirinya adalah $H$ dan $gH$, jadi $gH = G \setminus H$;
demikian pula $Hg = G \setminus H$. Karena itu $gH = Hg$ untuk setiap $g$: $H
\trianglelefteq G$.

(b) Menurut (a), $G/H$ adalah grup berorde $2$; kelas $\bar x$
memenuhi $\bar x^2 = \bar e$, yaitu $x^2 \in H$.

(c) Andaikan $H \leq A_4$ dengan $\abs H = 6$, sehingga berindeks $2$. Menurut
(b), $\sigma^2 \in H$ untuk setiap $\sigma \in A_4$. Setiap siklus-$3$
merupakan kuadrat semacam itu: jika $\sigma^3 = e$ maka $\sigma =
\sigma^4 = (\sigma^2)^2$. Jadi $H$ memuat kedelapan siklus-$3$ dari
$A_4$: $\abs H \geq 8 > 6$, sebuah kontradiksi. (Kebalikan teorema Lagrange
gagal pada kesempatan yang paling pertama: $6 \mid 12$.)
\end{solution}

\begin{solution}{exo:b3:groups:2}
Katakanlah $G/Z(G) = \langle gZ(G) \rangle$. Setiap $x \in G$ lalu ditulis
$x = g^k z$ dengan $k \in \Z$, $z \in Z(G)$. Untuk $x = g^kz$, $y =
g^l z'$:
\[
xy = g^k z\, g^l z' = g^{k+l} z z' = g^l z'\, g^k z = yx,
\]
unsur pusat komutatif dengan segalanya: jadi $G$ abelian.

Orde $p^2$: $Z(G) \neq \{e\}$ (\cref{thm:b3:groups:pfixed}), sehingga
$\abs{Z(G)} \in \{p, p^2\}$. Seandainya nilainya $p$, maka $G/Z(G)$ akan
berorde $p$, sehingga siklik, dan itu memaksa $G$ abelian serta $Z(G) =
G$ berorde $p^2$ --- kontradiksi. Jadi $Z(G) = G$.

Tak abelian berorde $p^3$: \emph{grup Heisenberg}
\[
H_p = \left\{ \begin{pmatrix} 1 & a & c\\ 0 & 1 & b\\ 0 & 0 & 1
\end{pmatrix} : a, b, c \in \mathbb F_p \right\}
\leq GL_3(\mathbb F_p),
\]
berorde $p^3$ (pemilihan $a, b, c$ bebas; ketertutupan dan inversnya lewat
perhitungan langsung). Grup ini tak abelian: kedua matriks elementer
$I + E_{12}$ dan $I + E_{23}$ mempunyai \omterm{def:b3:groups:derived}{komutator} $I + E_{13} \neq I$.
\end{solution}

\begin{solution}{exo:b3:groups:3}
(a) Sebuah morfisma $f \colon \Z/n\Z \to \Z/n\Z$ ditentukan oleh $k =
f(\bar 1)$ (lalu $f(\bar m) = m\bar k$), dan setiap $\bar k$ mendefinisikan
satu morfisma. Ia bijektif jika dan hanya jika $\bar k$ membangun $\Z/n\Z$, yaitu jika
$\gcd(k,n) = 1$, yaitu jika $\bar k \in (\Z/n\Z)^\times$. Komposisi
bersesuaian dengan perkalian: $f_k \circ f_l = f_{kl}$. Karena itu
$\operatorname{Aut}(\Z/n\Z) \cong (\Z/n\Z)^\times$.

(b) Pemetaan $\iota\colon G \to \operatorname{Aut}(G)$, $g \mapsto
\iota_g$, adalah morfisma: $\iota_g \circ \iota_h = \iota_{gh}$.
Petanya $\operatorname{Inn}(G)$; kernelnya $\{g : gxg^{-1} =
x\ \forall x\} = Z(G)$. Teorema isomorfisma pertama memberikan
$\operatorname{Inn}(G) \cong G/Z(G)$. Kenormalan di dalam
$\operatorname{Aut}(G)$: untuk $\alpha \in \operatorname{Aut}(G)$,
\[
(\alpha \circ \iota_g \circ \alpha^{-1})(x)
= \alpha\bigl(g\,\alpha^{-1}(x)\,g^{-1}\bigr)
= \alpha(g)\, x\, \alpha(g)^{-1} = \iota_{\alpha(g)}(x).
\]
\end{solution}

\begin{solution}{exo:b3:groups:4}
(a) Tinjau $\mu \colon H \times K \to HK$, $(h,k) \mapsto hk$, yang
surjektif menurut definisi. Tetapkan $h_0k_0 \in HK$: maka $hk = h_0k_0
\iff h_0^{-1}h = k_0 k^{-1} \in H \cap K$. Dengan menulis $u =
h_0^{-1}h$, serat dari $h_0k_0$ adalah $\{(h_0u,\, u^{-1}k_0) : u
\in H \cap K\}$, yang berkardinalitas $\abs{H \cap K}$. Karena itu $\abs
H\,\abs K = \abs{H\times K} = \abs{HK}\,\abs{H \cap K}$.

(b) Jika $HK$ subgrup: $KH \subseteq HK$ karena $kh =
\bigl(h^{-1}k^{-1}\bigr)^{-1} \in (HK)^{-1} = HK$; dan $HK
\subseteq KH$ dengan mengambil invers pada $HK = (HK)^{-1} \subseteq
(KH)^{-1}\dots$ lebih langsung lagi, untuk $hk \in HK$ berlaku $(hk)^{-1} =
k^{-1}h^{-1} \in KH$, jadi $HK = (HK)^{-1} \subseteq KH$; kedua
inklusi memberikan $HK = KH$. Sebaliknya jika $HK = KH$: ketertutupan,
$(hk)(h'k') = h(kh')k' \in h(HK)k' = (hH)(Kk') \subseteq HK$;
invers, $(hk)^{-1} = k^{-1}h^{-1} \in KH = HK$; dan $e \in HK$:
jadi subgrup. Jika misalnya $K \trianglelefteq G$, maka $hK = Kh$ untuk setiap
$h$, sehingga $HK = KH$ secara otomatis.

(c) Untuk $h \in H$, $k \in K$, \omterm{def:b3:groups:derived}{komutator} $[h,k] = hkh^{-1}
k^{-1}$ sama dengan $(hkh^{-1})k^{-1} \in K$ ($K$ \omterm{def:b3:groups:normal}{normal}) dan dengan
$h(kh^{-1}k^{-1}) \in H$ ($H$ \omterm{def:b3:groups:normal}{normal}), sehingga terletak di $H \cap K =
\{e\}$: jadi $hk = kh$.
\end{solution}

\begin{solution}{exo:b3:groups:5}
Cacahlah $E = \{(g,x) \in G \times X : g \cdot x = x\}$ dengan dua cara:
\[
\abs E = \sum_{g \in G} \abs{\operatorname{Fix}(g)}
= \sum_{x \in X} \abs{G_x}
= \sum_{x \in X} \frac{\abs G}{\abs{\mathcal O_x}}
= \abs G \sum_{\mathcal O \text{ orbit}} \sum_{x \in \mathcal O}
\frac{1}{\abs{\mathcal O}}
= \abs G \cdot \#\{\text{orbit}\},
\]
dengan memakai orbit--stabilisator ($\abs{G_x} = \abs G/\abs{\mathcal O_x}$)
dan pemartisian atas \omterm{def:b3:groups:action}{orbit}.

Kalung: misalkan $X$ himpunan semua pemetaan $\Z/p\Z \to
\{1, \dots, a\}$ (pewarnaan atas $p$ posisi), $\abs X = a^p$,
dengan $\Z/p\Z$ beraksi melalui rotasi. Identitasnya menetapkan seluruh $a^p$
pewarnaan. Rotasi $\bar k \neq \bar 0$ membangun $\Z/p\Z$ (sebab $p$
prima), jadi pewarnaan yang ditetapkannya invarian terhadap \emph{semua}
rotasi, sehingga konstan: ada $a$ pewarnaan tetap. Burnside memberikan
\[
\#\{\text{orbit}\} = \frac{a^p + (p-1)a}{p} \in \N ,
\]
jadi $p \mid a^p + (p-1)a$, yaitu $p \mid a^p - a$: teorema kecil
Fermat, murni lewat pencacahan.
\end{solution}

\begin{solution}{exo:b3:groups:6}
Orde $15 = 3 \cdot 5$: $n_3 \mid 5$ dan $n_3 \equiv 1 \pmod 3$
memaksa $n_3 = 1$; $n_5 \mid 3$ dan $n_5 \equiv 1 \pmod 5$ memaksa
$n_5 = 1$. Kedua \omterm{def:b3:groups:sylow}{subgrup Sylow} $P_3, P_5$ \omterm{def:b3:groups:normal}{normal}, berpotongan trivial
(ordenya relatif prima), dan $\abs{P_3P_5} = 15$
(\cref{exo:b3:groups:4}(a)): menurut \cref{prop:b3:groups:direct}, $G
\cong \Z/3\Z \times \Z/5\Z \cong \Z/15\Z$ (sisa Tiongkok).

Orde $45 = 3^2 \cdot 5$: $n_3 \mid 5$, $n_3 \equiv 1 \pmod 3$
memberikan $n_3 = 1$; $n_5 \mid 9$, $n_5 \equiv 1 \pmod 5$ memberikan $n_5 =
1$. Jadi $G \cong P_3 \times P_5$ dengan $\abs{P_3} = 9 = 3^2$ dan
$\abs{P_5} = 5$: keduanya abelian (\cref{thm:b3:groups:pfixed} untuk
$p^2$; orde prima berarti siklik), sehingga $G$ pun abelian.
\end{solution}

\begin{solution}{exo:b3:groups:7}
Orde $30 = 2 \cdot 3 \cdot 5$. $n_5 \mid 6$, $n_5 \equiv 1 \pmod
5$: $n_5 \in \{1, 6\}$; $n_3 \mid 10$, $n_3 \equiv 1 \pmod 3$:
$n_3 \in \{1, 10\}$. Andaikan $G$ \omterm{def:b3:groups:simple}{sederhana}, sehingga $n_5 = 6$ dan $n_3 =
10$. Dua subgrup berbeda yang berorde prima $p$ berpotongan trivial
(irisannya subgrup sejati dari $\Z/p\Z$), jadi keenam
subgrup Sylow-$5$ membawa $6 \times 4 = 24$ unsur berorde $5$,
dan kesepuluh subgrup Sylow-$3$ membawa $10 \times 2 = 20$ unsur
berorde $3$: $24 + 20 = 44 > 30 - 1$ unsur bukan identitas ---
mustahil. Jadi $n_5 = 1$ atau $n_3 = 1$: ada \omterm{def:b3:groups:sylow}{subgrup Sylow} yang \omterm{def:b3:groups:normal}{normal}.

Orde $56 = 2^3 \cdot 7$. $n_7 \mid 8$, $n_7 \equiv 1 \pmod 7$:
$n_7 \in \{1, 8\}$. Jika $n_7 = 8$, subgrup Sylow-$7$ membawa $8
\times 6 = 48$ unsur berorde $7$, sehingga tersisa tepat $56 - 48 =
8$ unsur lain. Subgrup Sylow-$2$ berorde $8$ dan terdiri
atas unsur-unsur itu, jadi ia \emph{sama dengan} himpunan tersebut: $n_2 = 1$.
Entah $n_7 = 1$ entah $n_2 = 1$: tak pernah \omterm{def:b3:groups:simple}{sederhana}.
\end{solution}

\begin{solution}{exo:b3:groups:8}
(a) \omterm{def:b3:groups:action}{Aksi} $g \cdot xH = gxH$ memberikan morfisma $\rho \colon G
\to \mathfrak S(G/H) \cong S_n$. Kernelnya adalah
\[
\ker\rho = \{g : \forall x \in G,\ gxH = xH\}
= \{g : \forall x,\ x^{-1}gx \in H\}
= \bigcap_{x \in G} xHx^{-1},
\]
yaitu \omterm{def:b3:groups:normal}{subgrup normal} (sebuah kernel) yang termuat di $H$ (ambil $x = e$). Jika
$N \trianglelefteq G$ dan $N \subseteq H$, maka untuk setiap $x$: $N
= xNx^{-1} \subseteq xHx^{-1}$, sehingga $N \subseteq \ker\rho$: jadi
kernel itulah yang terbesar.

(b) Misalkan $[G:H] = p$, pembagi prima terkecil dari $\abs G$, dan $K =
\ker\rho \subseteq H$. Maka $G/K$ terbenam di $S_p$, sehingga $[G:K]$
membagi $p!$. Selain itu $[G:K] = [G:H]\,[H:K] = p\,[H:K]$, jadi $[H:K]$
membagi $(p-1)!$. Padahal $[H:K]$ membagi $\abs G$, yang semua pembagi
primanya $\geq p$, sedangkan pembagi prima $(p-1)!$
semuanya $< p$: karena itu $[H:K] = 1$, yakni $H = K = \ker \rho$ bersifat
\omterm{def:b3:groups:normal}{normal}.

(c) Misalkan \omterm{def:b3:groups:action}{aksinya} transitif dan $x \in X$. Pemetaan $\Phi
\colon G/G_x \to X$, $gG_x \mapsto g \cdot x$, terdefinisi dengan baik dan
bijektif (orbit--stabilisator; \omterm{def:b3:groups:action}{orbitnya} seluruh $X$), dan ia
menganyam kedua \omterm{def:b3:groups:action}{aksi}: $\Phi(h \cdot gG_x) = \Phi(hgG_x) = (hg)
\cdot x = h \cdot \Phi(gG_x)$.
\end{solution}

\begin{solution}{exo:b3:groups:9}
(a) $D(G)$ \omterm{def:b3:groups:normal}{normal} dengan kuosien abelian
(\cref{prop:b3:groups:derived}); dan jika $N \trianglelefteq G$ mempunyai
$G/N$ abelian, proposisi yang sama memberikan $D(G) \subseteq N$: jadi
$D(G)$ yang terkecil. Sifat universalnya: misalkan $f \colon G \to A$
dengan $A$ abelian. Maka $f([x,y]) = [f(x), f(y)] = e$, sehingga $D(G)
\subseteq \ker f$, dan \cref{thm:b3:groups:firstiso} memfaktorkan $f =
\bar f \circ \pi$ melalui $G^{\mathrm{ab}}$, secara tunggal karena $\pi$
surjektif.

(b) \omterm{def:b3:groups:derived}{Komutator} adalah permutasi genap, jadi $D(S_n) \subseteq A_n$.
Sebaliknya setiap siklus-$3$ merupakan \omterm{def:b3:groups:derived}{komutator}:
\[
\bigl[(a\,b),\,(a\,c)\bigr] = (a\,b)(a\,c)(a\,b)(a\,c) =
(a\,b\,c),
\]
(periksa langsung pada $a, b, c$), dan siklus-$3$ membangun $A_n$
(\cref{lem:b3:groups:threecycles}): jadi $D(S_n) = A_n$ untuk $n \geq 3$,
dan $S_n^{\mathrm{ab}} \cong S_n/A_n \cong \Z/2\Z$. (Untuk $n = 2$:
$S_2$ abelian, $D(S_2) = \{e\}$, $S_2^{\mathrm{ab}} = S_2
\cong \Z/2\Z$ --- rumus $S_n^{\mathrm{ab}} \cong \Z/2\Z$
berlaku untuk setiap $n \geq 2$.)

$Q_8$: kuosien $Q_8/\{\pm 1\}$ berorde $4$, sehingga
abelian, jadi $D(Q_8) \subseteq \{\pm 1\}$; dan $[\mathrm i, \mathrm
j] = \mathrm i \mathrm j \mathrm i^{-1}\mathrm j^{-1} = \mathrm
i\mathrm j(-\mathrm i)(-\mathrm j) = (\mathrm i \mathrm j)^2 =
\mathrm k^2 = -1$, sehingga $D(Q_8) = \{\pm 1\}$ dan
$Q_8^{\mathrm{ab}} \cong (\Z/2\Z)^2$ (berorde $4$, bereksponen $2$:
kelas $\mathrm i, \mathrm j$ berkuadrat $\bar 1$).
\end{solution}

\begin{solution}{exo:b3:groups:10}
Induksi pada $\abs G$; untuk $\abs G = p$ satu-satunya subgrup
sejati adalah $H = \{e\}$, dan $N_G(\{e\}) = G \supsetneq
\{e\}$. Misalkan $Z = Z(G) \neq \{e\}$ (\cref{thm:b3:groups:pfixed}).

\emph{Jika $Z \not\subseteq H$}: pilih $z \in Z \setminus H$; unsur $z$
komutatif dengan $H$, jadi $zHz^{-1} = H$ dan $z \in N_G(H) \setminus
H$.

\emph{Jika $Z \subseteq H$}: beralihlah ke $\bar G = G/Z$, yaitu grup-$p$ yang
lebih kecil, dan $\bar H = H/Z \subsetneq \bar G$ (teorema
korespondensi). Dengan induksi, $N_{\bar G}(\bar H) \supsetneq \bar H$;
pilih $\bar g \in N_{\bar G}(\bar H) \setminus \bar H$ beserta angkatnya
$g$. Maka $g \notin H$, dan $gHg^{-1} \subseteq HZ = H$: memang
$\overline{ghg^{-1}} = \bar g \bar h \bar g^{-1} \in \bar H$ berarti
$ghg^{-1} \in HZ = H$ (sebab $Z \subseteq H$). Jadi $g \in
N_G(H)\setminus H$.

Subgrup maksimal: jika $M$ maksimal, $N_G(M) \supsetneq M$
memaksa $N_G(M) = G$: jadi $M \trianglelefteq G$. Maka $G/M$ adalah
grup-$p$ tanpa subgrup sejati yang tak trivial (korespondensi ditambah
kemaksimalan). Ambil $\bar x \neq \bar e$ di $G/M$, berorde $p^k$;
maka $\bar x^{p^{k-1}}$ membangun subgrup berorde $p$, yang
tentu seluruhnya: $\abs{G/M} = p$.
\end{solution}

\begin{solution}{exo:b3:groups:11}
(a) $\abs{A_5} = 60$. Tipe siklus di $A_5$: $e$; transposisi
ganda, ada $\frac{1}{2}\binom{5}{1}\binom{4}{2}\cdot 1 = 15$
buah ($5 \cdot 3$ cara: pilih titik tetapnya, lalu pasangkan sisanya);
siklus-$3$, ada $\frac{5 \cdot 4 \cdot 3}{3} = 20$; siklus-$5$, ada $4! =
24$.

Sebuah kelas $S_5$ yang termuat di $A_5$ entah tetap satu kelas-$A_5$
entah terpecah dua, bergantung pada apakah sentralisator-$S_5$ dari sebuah
unsur memuat permutasi ganjil ($\abs{\text{kelas di }A_5} =
60/\abs{Z_{A_5}(\sigma)}$ dan $Z_{A_5} = Z_{S_5} \cap A_5$).
Untuk $\sigma = (1\,2)(3\,4)$: $\abs{Z_{S_5}(\sigma)} = 120/15 = 8$,
dan $(1\,2) \in Z_{S_5}(\sigma)$ ganjil, jadi $\abs{Z_{A_5}} = 4$
dan kelasnya beranggota $60/4 = 15$ unsur: tanpa pemecahan. Untuk $\sigma =
(1\,2\,3)$: $Z_{S_5}(\sigma) \supseteq \langle \sigma \rangle
\times \langle (4\,5)\rangle$, berorde $6 = 120/20$, sehingga keduanya sama;
kelompok ini memuat $(4\,5)$ yang ganjil: kelasnya $60/3 = 20$: tanpa pemecahan. Untuk
$\sigma$ berupa siklus-$5$: $Z_{S_5}(\sigma) = \langle \sigma\rangle$
(berorde $120/24 = 5$), semuanya genap: $Z_{A_5}(\sigma) = \langle
\sigma\rangle$ dan kelas-$A_5$ beranggota $60/5 = 12$ unsur --- jadi
$24$ siklus lima terpecah menjadi \emph{dua} kelas beranggota $12$. Ukuran
kelasnya: $1, 15, 20, 12, 12$.

(b) \omterm{def:b3:groups:normal}{Subgrup normal} $N$ adalah gabungan \omterm{ex:b3:groups:actions}{kelas konjugasi}
yang memuat $\{e\}$, dengan $\abs N \mid 60$. Jumlah yang mungkin, $1 +
(\text{himpunan bagian dari } \{15, 20, 12, 12\})$, adalah
\[
1,\ 13,\ 13,\ 16,\ 21,\ 25,\ 28,\ 28,\ 33,\ 36,\ 40,\ 40,\ 45,\
48,\ 48,\ 60 ;
\]
dan satu-satunya pembagi $60$ dalam daftar itu adalah $1$ dan $60$: jadi $N =
\{e\}$ atau $A_5$.

(c) Misalkan $G$ \omterm{def:b3:groups:simple}{sederhana} dengan $\abs G = 60$. $n_5 \mid 12$, $n_5
\equiv 1 \pmod 5$: $n_5 \in \{1, 6\}$. Nilai $n_5 = 1$ akan membuat
subgrup Sylow-$5$ \omterm{def:b3:groups:normal}{normal}, dan itu bertentangan dengan kesederhanaan ($1 < 5 <
60$). Karena itu $n_5 = 6$.
\end{solution}

\begin{solution}{exo:b3:groups:12}
(a) $P$ \omterm{def:b3:groups:normal}{normal} di $H = N_G(P)$ menurut definisi
\omterm{ex:b3:groups:actions}{normalisator}, dan ia subgrup Sylow-$p$ dari $H$ (ordenya sudah
merupakan seluruh bagian-$p$ dari $\abs G$, apalagi dari $\abs H$).
\omterm{def:b3:groups:sylow}{Subgrup Sylow} yang \omterm{def:b3:groups:normal}{normal} bersifat tunggal: yang lain akan
konjugat dengannya (Sylow II di dalam $H$), sehingga sama dengannya.

(b) Misalkan $g \in N_G(H)$. Maka $gPg^{-1} \subseteq gHg^{-1} = H$
adalah subgrup $H$ yang berorde sama dengan $P$: jadi subgrup
Sylow-$p$ dari $H$, sehingga $gPg^{-1} = P$ menurut (a). Dengan demikian $g \in
N_G(P) = H$: $N_G(H) \subseteq H$, dan inklusi sebaliknya
trivial.

(c) Andaikan $H \subseteq N \trianglelefteq G$ dengan $N$ sejati.
Di sini $P$ subgrup Sylow-$p$ dari $N$; untuk sembarang $g \in G$,
$gPg^{-1} \subseteq N$ juga demikian, jadi $gPg^{-1} = nPn^{-1}$ untuk
suatu $n \in N$ (Sylow II di dalam $N$), sehingga $n^{-1}g \in N_G(P)
\subseteq N$ dan $g \in N$: jadi $N = G$, kontradiksi (inilah
\emph{argumen Frattini}). Untuk subgrup maksimal $M \supseteq
N_G(P)$: $N_G(M) \supseteq M$ sama dengan $M$ atau $G$; jika $G$,
maka $M \trianglelefteq G$ adalah \omterm{def:b3:groups:normal}{subgrup normal} sejati yang
memuat $N_G(P)$ --- dan itu sudah dikecualikan pada butir sebelumnya. Jadi
$N_G(M) = M$.
\end{solution}

\begin{solution}{pb:b3:groups:1}
\textbf{1.} Untuk $x, y \in G$: $(xy)^2 = e$ memberikan $xy = (xy)^{-1} =
y^{-1}x^{-1} = yx$ (setiap unsur menjadi inversnya sendiri): jadi abelian.
Grup $G$ semacam itu, jika ditulis secara aditif, adalah ruang vektor atas $\mathbb
F_2$ ($2x = 0$, dan aksiomanya persis aksioma grup abelian); jika
berhingga, ia mempunyai basis berhingga: $G \cong (\Z/2\Z)^k$, berorde
$2^k$.

\textbf{2.} Orde $p^2$: $G$ abelian
(\cref{thm:b3:groups:pfixed}). Jika ada unsur berorde $p^2$, maka
$G$ siklik. Jika tidak, semua $x \neq e$ berorde $p$; secara aditif
$G$ lalu menjadi ruang vektor atas $\mathbb F_p$ ($px = 0$), berdimensi
$2$ (karena ada $p^2$ unsur): $G \cong (\Z/p\Z)^2$.

Yang abelian berorde $8$, diuraikan menurut orde terbesar $m$ sebuah unsur: $m =
8$: siklik $\Z/8\Z$. $m = 2$: $(\Z/2\Z)^3$ menurut pertanyaan 1. $m = 4$:
misalkan $x$ berorde $4$ dan $y \notin \langle x \rangle$; maka $y^2 \in
\langle x\rangle$ (indeks $2$). Nilai $y^2 \in \{x, x^3\}$ akan membuat $y$
berorde $8$; jadi $y^2 \in \{e, x^2\}$. Jika $y^2 = x^2$, gantilah $y$ dengan
$xy$: $(xy)^2 = x^2y^2 = x^4 = e$ (sebab $G$ abelian) dan $xy \notin
\langle x\rangle$. Jadi kita boleh menganggap $y^2 = e$: maka $\langle x
\rangle \cap \langle y \rangle = \{e\}$, keduanya \omterm{def:b3:groups:normal}{normal} (karena abelian),
$\abs{\langle x\rangle \langle y\rangle} = 8$:
\cref{prop:b3:groups:direct} memberikan $G \cong \Z/4\Z \times \Z/2\Z$.
Tanpa pengulangan: banyaknya penyelesaian $x^2 = e$ berturut-turut $2, 4, 8$
pada ketiga grup itu.

\textbf{3.} Definisikan $\psi \colon H \rtimes_{\varphi'} K \to H
\rtimes_{\varphi} K$ dengan $\psi(h, k) = (h, \alpha(k))$, sebuah bijeksi.
Ia morfisma:
\[
\psi\bigl((h,k)(h',k')\bigr)
= \bigl(h\,\varphi'(k)(h'),\, \alpha(kk')\bigr)
= \bigl(h\,\varphi(\alpha k)(h'),\, \alpha(k)\alpha(k')\bigr)
= \psi(h,k)\,\psi(h',k').
\]

\textbf{4.} $\operatorname{Aut}(\Z/n\Z) \cong (\Z/n\Z)^\times$
(\cref{exo:b3:groups:3}): untuk $n = 3$: $\{\pm 1\} \cong \Z/2\Z$;
$n = 4$: $\{\bar 1, \bar 3\} \cong \Z/2\Z$; $n = 5$: $\{\bar 1,
\bar 2, \bar 3, \bar 4\}$, siklik berorde $4$ yang dibangun oleh $\bar
2$ ($2, 4, 3, 1$); $n = 7$: siklik berorde $6$ yang dibangun oleh $\bar
3$ ($3, 2, 6, 4, 5, 1$). Untuk $V = (\Z/2\Z)^2$: automorfismanya bersifat
$\mathbb F_2$-linear (ia mempertahankan penjumlahan, dan skalarnya hanya $0,
1$), sehingga $\operatorname{Aut}(V) = GL_2(\mathbb F_2)$, berorde
$(4-1)(4-2) = 6$; ia beraksi secara setia pada $3$ vektor tak nol,
sehingga memberikan morfisma injektif ke $S_3$ antara dua grup berorde
$6$: jadi $\operatorname{Aut}(V) \cong S_3$.

\textbf{5.} Teorema Cauchy menyediakan $r$ berorde $p$; $N = \langle r
\rangle$ berindeks $2$, sehingga \omterm{def:b3:groups:normal}{normal} (\cref{exo:b3:groups:1}).
Cauchy juga menyediakan $s$ berorde $2$, dan $s \notin N$ (semua
unsur $N$ selain identitas berorde ganjil $p$).

\textbf{6.} $N \cap \langle s \rangle = \{e\}$ dan $\abs{N\langle
s\rangle} = 2p$ (\cref{exo:b3:groups:4}(a)): menurut
\cref{prop:b3:groups:semidirect}, $G \cong \Z/p\Z \rtimes_\varphi
\Z/2\Z$ dengan $\varphi(1) = (x \mapsto sxs^{-1})$ sebuah automorfisma
yang ordenya membagi $2$. Di $(\Z/p\Z)^\times$, persamaan $k^2 = 1$ hanya mempunyai
penyelesaian $k = \pm 1$ (sebab $X^2 - 1$ berakar paling banyak dua di
lapangan $\mathbb F_p$). Jika $\varphi(1) = \mathrm{id}$: hasil kalinya
langsung, $G \cong \Z/p\Z \times \Z/2\Z \cong \Z/2p\Z$. Jika
$\varphi(1) = -\mathrm{id}$: $G = \langle r, s \mid r^p = s^2 = e,
\ srs^{-1} = r^{-1}\rangle \cong D_p$
(\cref{ex:b3:groups:semidirectexamples}). Keduanya tidak isomorfik:
$D_p$ tak abelian untuk $p \geq 3$ ($srs^{-1} = r^{-1} \neq r$).

\textbf{7.} Bilangan $n_q \equiv 1 \pmod q$ membagi $p < q$: jadi $n_q = 1$, sehingga
$N \cong \Z/q\Z$ \omterm{def:b3:groups:normal}{normal}. Misalkan $P \cong \Z/p\Z$ sebuah subgrup
Sylow-$p$: $N \cap P = \{e\}$, $NP = G$ (berorde $pq$), jadi $G
\cong \Z/q\Z \rtimes_\varphi \Z/p\Z$ dengan $\varphi \colon \Z/p\Z
\to \operatorname{Aut}(\Z/q\Z) \cong \Z/(q-1)\Z$ (siklik,
diterima tanpa bukti). Jika $p \nmid q - 1$: peta $\varphi$ berorde
pembagi $p$ sekaligus $q-1$, sehingga trivial, dan $G \cong
\Z/pq\Z$ (\cref{ex:b3:groups:pq}). Jika $p \mid q - 1$: selain
$\varphi$ yang trivial, setiap $\varphi$ tak trivial bersifat injektif (sebab
kernelnya, subgrup $\Z/p\Z$, trivial) dengan peta berupa
satu-satunya subgrup $C$ berorde $p$ dari grup siklik
$\Z/(q-1)\Z$. Dua \omterm{def:b3:groups:action}{aksi} tak trivial $\varphi, \varphi'$ lalu berupa
dua isomorfisma $\Z/p\Z \to C$, sehingga $\alpha =
\varphi^{-1}\circ\varphi' \in \operatorname{Aut}(\Z/p\Z)$
memenuhi $\varphi' = \varphi\circ\alpha$: menurut pertanyaan 3, kedua
\omterm{def:b3:groups:semidirect}{hasil kali semilangsung} itu isomorfik. Karena itu ada tepat satu grup tak abelian
berorde $pq$ (tak abelian sebab $\varphi \neq
\mathrm{id}$ membuat suatu konjugasi menjadi tak trivial). Orde $15$: $p =
3$, $q = 5$, $3 \nmid 4$: hanya yang siklik.

\textbf{8.} Tidak setiap unsur berorde $\leq 2$ (kalau demikian ia abelian menurut
pertanyaan 1), dan tidak ada unsur berorde $8$ (kalau ada ia siklik, jadi abelian):
jadi ada $r$ berorde $4$, dan $N = \langle r\rangle$, yang berindeks $2$,
bersifat \omterm{def:b3:groups:normal}{normal}.

\textbf{9.} $s^2 \in N$ menurut \cref{exo:b3:groups:1}(b). Konjugat
$srs^{-1} \in N$ berorde $4$, jadi $srs^{-1} \in \{r,
r^3\}$; jika $srs^{-1} = r$ maka $r$ dan $s$ komutatif sehingga $G =
\langle r, s\rangle$ abelian --- dan itu dikecualikan. Jadi $srs^{-1} =
r^{-1}$. Jika $s^2 = r$ atau $r^3$, maka $s$ berorde $8$: dikecualikan.
(Cara lain: $s^2$ komutatif dengan $s$, padahal $s r s^{-1} =
r^{-1}$ dan $s r^3 s^{-1} = r^{-3} = r$: baik $r$ maupun $r^3$ tidak
ditetapkan oleh konjugasi dengan $s$.) Jadi $s^2 \in \{e, r^2\}$.

\textbf{10.} Jika $s^2 = e$: $G = \langle r, s \mid r^4 = s^2 = e,\
srs^{-1} = r^{-1}\rangle$. Kedelapan unsur $r^is^j$ ($0 \leq i
< 4$, $0 \leq j < 2$) berbeda ($s \notin \langle r\rangle$)
dan relasinya menentukan seluruh hasil kali: pengaitan $r
\mapsto$ (rotasi sebesar $\pi/2$), $s \mapsto$ (sebuah pencerminan) mendefinisikan
morfisma surjektif pada $D_4$, antara dua grup berorde $8$: jadi sebuah
isomorfisma.

\textbf{11.} Jika $s^2 = r^2$: kembali $G = \{r^i s^j\}$ dan
relasi $r^4 = e$, $s^2 = r^2$, $srs^{-1} = r^{-1}$ memaksa
seluruh tabelnya. Dengan $\mathrm i = r$, $\mathrm j = s$, $\mathrm k =
rs$, $-1 = r^2$: $\mathrm i^2 = \mathrm j^2 = -1$, $\mathrm k^2 =
rsrs = r r^{-1} s s = s^2 = -1$ (memakai $sr = r^{-1}s$), dan
$\mathrm i \mathrm j \mathrm k = r\,s\,rs = r\,r^{-1}s\,s = s^2 =
-1$. Keberadaannya: matriks
\[
A = \begin{pmatrix} \iu & 0\\ 0 & -\iu \end{pmatrix},
\qquad
B = \begin{pmatrix} 0 & 1\\ -1 & 0 \end{pmatrix}
\]
memenuhi $A^4 = I$, $B^2 = -I = A^2$ dan $BAB^{-1} = A^{-1}$ ---
untuk yang terakhir, periksalah
\[
BA = \begin{pmatrix} 0 & -\iu\\ -\iu & 0 \end{pmatrix} = A^{-1}B .
\]
Jadi $\{\pm I, \pm A, \pm B, \pm AB\}$ adalah grup berorde $8$ yang
mewujudkan tabel itu: $Q_8$ memang ada.

\textbf{12.} Misalkan $H \neq \{e\}$ sebuah subgrup dan $x \in H
\setminus \{e\}$. Jika $x \neq -1$ maka $x \in \{\pm\mathrm i,
\pm\mathrm j, \pm\mathrm k\}$ dan $x^2 = -1 \in H$. Jadi $-1 \in H$
selalu. Subgrupnya adalah $\{e\}$, $\{\pm 1\}$ (pusatnya),
$\langle \mathrm i\rangle, \langle \mathrm j\rangle, \langle
\mathrm k\rangle$ (berindeks $2$) dan $Q_8$: semuanya \omterm{def:b3:groups:normal}{normal} ($\{e\}$ dan
pusatnya secara trivial, yang berindeks $2$ menurut \cref{exo:b3:groups:1}, dan $Q_8$
sendiri). Grup $D_4$ mempunyai lima unsur berorde $2$ ($r^2$ dan keempat
pencerminan), sedangkan $Q_8$ hanya satu ($-1$): jadi tak isomorfik. \omterm{def:b3:groups:semidirect}{Hasil kali
semilangsung} $H \rtimes K$ dengan $H, K \neq \{e\}$ menuntut $H \cap K =
\{e\}$, dan itu mustahil karena keduanya memuat $-1$.

\textbf{13.} $n_3 \mid 4$, $n_3 \equiv 1 \pmod 3$: $n_3 \in \{1,
4\}$; $n_2 \mid 3$, ganjil: $n_2 \in \{1, 3\}$. Jika $n_3 = 4$: keempat
subgrup Sylow-$3$ berpotongan trivial berpasangan (karena berorde
prima), sehingga memberikan $4 \times 2 = 8$ unsur berorde $3$; sisa
$4$ unsur pastilah membentuk satu-satunya subgrup
Sylow-$2$: jadi $n_2 = 1$.

\textbf{14.} Kernel $\ker\rho$ menormalkan setiap subgrup Sylow-$3$, jadi
$\ker \rho \subseteq N_G(P_3)$, yang berindeks $n_3 = 4$, yakni
berorde $3$: $\abs{\ker\rho} \in \{1, 3\}$. Orde $3$ akan membuat
$\ker\rho$ menjadi subgrup Sylow-$3$ yang \emph{\omterm{def:b3:groups:normal}{normal}}, dan itu bertentangan dengan $n_3
= 4$. Jadi $\rho$ injektif dan petanya $H \leq S_4$ berorde
$12$, berindeks $2$: $H \trianglelefteq S_4$ dan $H$ memuat semua
kuadrat (\cref{exo:b3:groups:1}(b)). Kuadrat di $S_4$ mencakup
$e$ dan kedelapan siklus-$3$ (sebab $\sigma = (\sigma^2)^2$ untuk sebuah
siklus-$3$), yang membangun $A_4$ (mereka terletak di $A_4$, dan bersama
hasil kalinya memberikan seluruh dua belas unsur; atau:
\cref{lem:b3:groups:threecycles} pada bagian pembangkitan untuk $n = 4$,
yang hanya memerlukan $n \geq 3$). Jadi $A_4 \subseteq H$ dan $\abs{A_4}
= \abs H$: $G \cong H = A_4$.

\textbf{15.} $P_3 \trianglelefteq G$, $P_3 \cap P_2 = \{e\}$,
$P_3P_2 = G$: jadi $G \cong \Z/3\Z \rtimes_\varphi P_2$
(\cref{prop:b3:groups:semidirect}), dengan $\varphi \colon P_2 \to
\operatorname{Aut}(\Z/3\Z) = \{\pm\mathrm{id}\} \cong \Z/2\Z$.
\begin{itemize}
  \item $\varphi$ trivial: \omterm{prop:b3:groups:direct}{hasil kali langsung} $\Z/3\Z \times \Z/4\Z
        \cong \Z/12\Z$ dan $\Z/3\Z \times (\Z/2\Z)^2 \cong \Z/6\Z
        \times \Z/2\Z$.
  \item $P_2 = \Z/4\Z$, $\varphi$ surjektif: mau tak mau
        $\varphi(1) = -\mathrm{id}$ (satu-satunya pilihan tak trivial):
        satu grup, yaitu $\mathrm{Dic}_3 = \Z/3\Z \rtimes \Z/4\Z$.
  \item $P_2 = (\Z/2\Z)^2$, $\varphi$ surjektif: di sini $\ker\varphi$
        salah satu dari tiga subgrup berorde $2$; ketiga
        $\varphi$ yang dihasilkan berbeda oleh automorfisma $(\Z/2\Z)^2$
        yang mempermutasikan ketiga subgrup itu (pertanyaan 4: $\operatorname{Aut}
        \cong S_3$ beraksi transitif pada ketiga involusinya), sehingga
        menurut pertanyaan 3 semuanya memberi satu kelas isomorfisma. Kelas itu
        $D_6$: pilih $t$ yang membangun $\ker\varphi$ dan $x$
        yang membangun $\Z/3\Z$; unsur $\rho = (x, t)$ memenuhi
        $\rho^2 = (2x, 0)$, $\rho^3 = (0, t)$, $\rho^6 = e$ dan tidak ada
        pangkat lebih kecil yang sama dengan $e$: jadi berorde $6$; untuk $s = (0, u)$ dengan $u
        \notin \ker\varphi$: $s^2 = e$ dan $s\rho s^{-1} = (-x, t)
        = \rho^{-1}$. Karena $\langle \rho, s\rangle$ berorde $12$,
        maka $G \cong D_6$.
\end{itemize}

\textbf{16.} Dengan mencacah unsur berorde $2$: $\Z/12\Z$ mempunyai $1$;
$\Z/6\Z\times\Z/2\Z$ mempunyai $3$; $D_6$ mempunyai $7$ (enam pencerminan dan
setengah putaran $\rho^3$); $A_4$ mempunyai $3$; $\mathrm{Dic}_3$ mempunyai $1$
(hanya $(0, 2)$: unsur $(h, k)$ dengan $k$ berorde $4$ di
$\Z/4\Z$ berorde $4$). Ini memisahkan semuanya kecuali pasangan
$\{\Z/12\Z, \mathrm{Dic}_3\}$ dan $\{\Z/6\Z\times\Z/2\Z, A_4\}$:
anggota pertamanya abelian, yang kedua tidak ($\mathrm{Dic}_3$:
\omterm{def:b3:groups:action}{aksinya} tak trivial; $A_4$: $(1\,2\,3)$ dan
$(1\,2)(3\,4)$ tidak komutatif). Jadi lima grup berbeda; Bagian
II--IV menunjukkan bahwa daftarnya lengkap.

\textbf{17.} Tabel klasifikasinya:
\begin{center}
\begin{tabular}{c l c}
$n$ & grup berorde $n$ & \#\\
\hline
$1$ & $\{e\}$ & $1$\\
$2$ & $\Z/2\Z$ & $1$\\
$3$ & $\Z/3\Z$ & $1$\\
$4$ & $\Z/4\Z$, $(\Z/2\Z)^2$ & $2$\\
$5$ & $\Z/5\Z$ & $1$\\
$6$ & $\Z/6\Z$, $S_3 = D_3$ & $2$\\
$7$ & $\Z/7\Z$ & $1$\\
$8$ & $\Z/8\Z$, $\Z/4\Z{\times}\Z/2\Z$, $(\Z/2\Z)^3$, $D_4$,
      $Q_8$ & $5$\\
$9$ & $\Z/9\Z$, $(\Z/3\Z)^2$ & $2$\\
$10$ & $\Z/10\Z$, $D_5$ & $2$\\
$11$ & $\Z/11\Z$ & $1$\\
$12$ & $\Z/12\Z$, $\Z/6\Z{\times}\Z/2\Z$, $D_6$, $A_4$,
       $\mathrm{Dic}_3$ & $5$\\
$13$ & $\Z/13\Z$ & $1$\\
$14$ & $\Z/14\Z$, $D_7$ & $2$\\
$15$ & $\Z/15\Z$ & $1$\\
\end{tabular}
\end{center}
Orde $6, 10, 14$ berasal dari Bagian II dengan $p = 3, 5, 7$; orde $15$ dari
pertanyaan 7; orde $8$ dari Bagian III bersama pertanyaan 2; orde
$12$ dari Bagian IV; orde prima dari Lagrange; orde $4$ dan $9$
dari pertanyaan 2.

\textbf{18.} Pusat $Z = Z(G)$ tak trivial
(\cref{thm:b3:groups:pfixed}) dan $Z \neq G$ (karena tak abelian), jadi
$\abs Z \in \{p, p^2\}$. Jika $\abs Z = p^2$, maka $G/Z$ siklik
berorde $p$ dan \cref{exo:b3:groups:2} membuat $G$ abelian:
dikecualikan, sehingga $\abs Z = p$ dan $\abs{G/Z} = p^2$. Menurut pertanyaan 2,
$G/Z$ berupa $\Z/p^2\Z$ atau $(\Z/p\Z)^2$; yang siklik sekali lagi dikecualikan oleh
\cref{exo:b3:groups:2}: jadi $G/Z \cong (\Z/p\Z)^2$. Karena $G/Z$
abelian, setiap \omterm{def:b3:groups:derived}{komutator} terletak di $Z$: $D(G) \subseteq Z$; dan
$D(G) \neq \{e\}$ (sebab $G$ tak abelian), sehingga $D(G) = Z$ (nilai $\abs Z = p$
tidak menyisakan ruang lain). \omterm{def:b3:groups:derived}{Komutatornya} bersifat pusat dengan orde membagi
$\abs Z = p$.

\textbf{19.} Induksi pada $k$, dengan kasus $k = 1$ yang trivial.
Dengan memakai $yx = xyz^{-1}\cdot$\,--- lebih tepatnya, $z = [y, x] =
yxy^{-1}x^{-1}$ memberikan $yx = zxy$, yakni memindahkan satu $y$ ke kiri
melewati satu $x$ menghasilkan satu faktor $z$, yang bersifat pusat sehingga dapat
diparkir di mana saja. Maka
\[
(xy)^{k+1} = (xy)^k\,xy = x^ky^kz^{k(k-1)/2}\,xy
= x^k\,(y^kx)\,y\,z^{k(k-1)/2}
= x^{k+1}y^{k+1}\,z^{k(k-1)/2 + k},
\]
sebab membawa $x$ melewati $y^k$ berharga $k$ faktor $z$ ($y^kx =
z^kxy^k$, lewat $k$ penerapan $yx = zxy$); dan $k(k-1)/2 + k
= k(k+1)/2$.

\textbf{20.} Dengan $k = p$: $(xy)^p = x^py^pz^{p(p-1)/2}$. Untuk
$p$ ganjil, $(p-1)/2$ bilangan bulat, sehingga $z^{p(p-1)/2} =
(z^p)^{(p-1)/2} = e$ (pertanyaan 18: orde $z$ membagi $p$):
jadi $\theta(xy) = \theta(x)\theta(y)$, sebuah morfisma. Nilainya bersifat
pusat: kelas $x$ di $G/Z \cong (\Z/p\Z)^2$ berorde
pembagi $p$, sehingga $x^p \in Z$. Jika $\theta$ trivial, setiap
unsur berorde pembagi $p$: jadi eksponennya $p$ (bukan $1$: sebab $G \neq
\{e\}$). Jika tidak, ada $x^p \neq e$, dan $x$ berorde $p^2$
(ordenya membagi $p^3$, dan $x$ tidak mungkin berorde $p^3$: sebab $G$
akan siklik, jadi abelian): eksponennya $p^2$.

\textbf{21.} Kelas $\bar x, \bar y$ membangun $G/Z$: \omterm{def:b3:groups:derived}{komutator}
keduanya $z = [y, x]$ tak trivial, yakni $\neq e$, sebab kalau tidak $x, y, Z$ akan
membangun $G$ yang abelian (kelasnya membangun \omterm{thm:b3:groups:quotient}{grup kuosiennya} dan
$Z$ bersifat pusat) --- lagi pula $z$ membangun $Z$ (karena $\abs Z = p$). Setiap
$g \in G$ berkelas $\bar x^a\bar y^b$ untuk $0 \leq a, b <
p$ yang tunggal, sehingga $g = x^ay^bz^c$ dengan $0 \leq c < p$ yang tunggal: jadi $p^3$
unsur, semuanya terhitung. Hasil kali bentuk \omterm{def:b3:groups:normal}{normal} semacam itu
dihitung hanya dengan $yx = zxy$, $z$ bersifat pusat, dan $x^p = y^p =
z^p = e$: tabelnya terpaksa demikian, sehingga dua grup tak abelian
bereksponen-$p$ berorde $p^3$ pastilah isomorfik (cocokkan
pembangunnya). Grup Heisenberg mewujudkan relasi itu: dengan
$X = I + E_{12}$, $Y = I + E_{23}$, diperoleh $[Y, X] = I -
E_{13}$ (yang bersifat pusat di $H_p$), dan untuk sembarang $N$ segitiga atas
sejati, $N^3 = 0$ memberikan
\[
(I + N)^p = I + pN + \binom p2N^2 = I
\qquad\text{pada karakteristik } p,\ p \geq 3,
\]
sebab $p \mid p$ dan $p \mid \binom p2$ untuk $p$ ganjil: jadi eksponennya
$p$. Karena itu grup bereksponen-$p$ tersebut adalah $H_p$.

\textbf{22.} $Z(G) = \langle r^p\rangle$: memang $r^p$ bersifat
pusat (lewat argumen pertanyaan 20: kelas $r$ di
\omterm{thm:b3:groups:quotient}{grup kuosien} bereksponen-$p$ $G/Z$ memberikan $r^p \in Z$) dan bernilai $\neq e$,
sehingga ia membangun pusat yang berorde $p$. Ambil sembarang $t \notin N$. Jika
$t^p = e$, tetapkan $s = t$. Jika tidak, $\theta(t) = t^p \in Z =
\langle r^p\rangle$, katakanlah $t^p = r^{pb}$; tetapkan $s = tr^{-b}$: menurut
pertanyaan 20 ($\theta$ morfisma, $p$ ganjil), $s^p =
t^pr^{-pb} = e$, dan $s \notin N$. Konjugasinya: $srs^{-1} \in
N$ ($N$ \omterm{def:b3:groups:normal}{normal}) berorde $p^2$, jadi $srs^{-1} = r^m$ dengan $p
\nmid m$; lagi pula $s^p = e$ memaksa $m^p \equiv m \pmod{p^2}$ ---
mengonjugasikan $p$ kali mengembalikan $r$, sehingga $m^p \equiv 1 \pmod
{p^2}$, dan $m \equiv m^p \equiv 1 \pmod p$ (Fermat): jadi $m = 1 +
ap$. Ketaktrivialan (sebab $G$ tak abelian) memberikan $a \not\equiv 0$;
mengganti $s$ dengan pangkat $s^{a'}$ dengan $aa' \equiv 1 \pmod p$
mengubah \omterm{def:b3:groups:action}{aksinya} menjadi $r \mapsto r^{1+p}$. Ini menyajikan $G$ sebagai
$\Z/p^2\Z\rtimes_\varphi\Z/p\Z$ dengan $\varphi(1)\colon r
\mapsto r^{1+p}$; menurut pertanyaan 3, dua morfisma tak trivial
$\Z/p\Z \to \operatorname{Aut}(\Z/p^2\Z)$ dengan peta yang sama
--- dan petanya \emph{satu-satunya} subgrup berorde $p$ dari
$\operatorname{Aut}(\Z/p^2\Z)$ yang siklik --- memberikan \omterm{def:b3:groups:semidirect}{hasil kali
semilangsung} yang isomorfik: jadi tunggal.

\textbf{23.} Yang abelian: $\Z/p^3\Z$, $\Z/p^2\Z\times\Z/p\Z$,
$(\Z/p\Z)^3$ (argumen pertanyaan 2, naik satu derajat: uraikan menurut
orde maksimal). Yang tak abelian: tepat $H_p$ (bereksponen $p$,
pertanyaan 21) dan $\Z/p^2\Z\rtimes\Z/p\Z$ (bereksponen $p^2$,
pertanyaan 22), yang dibedakan oleh eksponennya. Totalnya: lima.
Untuk $p = 2$ argumen morfisma pada pertanyaan 20 runtuh:
$z^{p(p-1)/2} = z^{1} = z \neq e$, pengkuadratan bukan morfisma,
dan memang kedua grup tak abelian berorde $8$ bereksponen
$4$ --- invarian yang memisahkan $D_4$ dari $Q_8$ adalah
banyaknya unsur berorde $2$ (lima lawan satu), bukan
eksponennya. Dunia $p$ ganjil, untuk sekali ini, lebih rapi daripada
karakteristik $2$.

\textbf{24.} Di $H_p$ setiap unsur $\neq e$ berorde $p$
(karena bereksponen $p$, pertanyaan 21): jadi ada $p^3 - 1$ unsur berorde $p$.
Di $M_p$, pemetaan $\theta \colon x \mapsto x^p$ adalah morfisma
$M_p \to Z(M_p) = \langle r^p\rangle$ (pertanyaan 20, $p$ ganjil);
$\theta(r) = r^p \neq e$, sehingga petanya seluruh pusat yang berorde $p$
dan $\ker\theta = \{x : x^p = e\}$ berorde $p^3/p =
p^2$. Unsur berorde $p$ adalah unsur bukan identitas
dari kernel itu: jadi ada $p^2 - 1$ buah. Untuk $p = 3$: $H_3$ mempunyai
$27 - 1 = 26$ unsur berorde $3$, dan $M_3 = \Z/9\Z \rtimes
\Z/3\Z$ mempunyai $9 - 1 = 8$. Untuk $p = 2$ argumennya mati sejak
awal: pengkuadratan bukan morfisma pada grup tak abelian berorde
$8$ (pertanyaan 23), dan memang himpunan $\{x : x^2 = e\}$
beranggota $6$ unsur di $D_4$ --- bukan orde sebuah subgrup
$D_4$. Cacahan yang memang memisahkan pasangan itu adalah banyaknya
unsur berorde $2$: lima di $D_4$, satu di $Q_8$ (pertanyaan
11).

\textbf{25.} Jika $p^2 \mid n$, grup $\Z/p\Z \times
\Z/(n/p)\Z$ berorde $n$ dan tidak siklik: orde setiap unsurnya
membagi $\operatorname{lcm}(p, n/p) = n/p < n$, sebab
$p \mid n/p$. Jika $p < q$ bilangan prima yang membagi $n$ dengan $p \mid
q - 1$, pertanyaan 7 menyediakan grup tak abelian $\Z/q\Z \rtimes
\Z/p\Z$ berorde $pq$; maka $(\Z/q\Z \rtimes \Z/p\Z) \times
\Z/(n/pq)\Z$ berorde $n$ dan tak abelian, sehingga tidak
siklik. Sekarang andaikan $\gcd(n, \varphi(n)) > 1$ lalu pilih bilangan prima
$p$ yang membagi keduanya. Dengan menulis $\varphi(n) = \prod_{q^a \parallel
n} q^{a-1}(q - 1)$, keterbagian $p \mid \varphi(n)$
berarti entah $p^2 \mid n$ (dari faktor $q^{a-1}$ dengan $q = p$,
$a \geq 2$) entah $p \mid q - 1$ untuk suatu bilangan prima $q \mid n$, $q
\neq p$: pada kedua kasus $n$ tidak siklik menurut uraian di atas. Lewat
kontraposisi, $n$ yang siklik memaksa $\gcd(n, \varphi(n)) = 1$.
Periksa untuk $n \leq 15$: nilai $\varphi(n)$ untuk $n = 1,
\dots, 15$ adalah $1, 1, 2, 2, 4, 2, 6, 4, 6, 4, 10, 4, 12, 6,
8$, dan $\gcd(n, \varphi(n)) = 1$ tepat untuk $n = 1, 2, 3, 5,
7, 11, 13, 15$ --- persis entri tabel pada
pertanyaan 17 yang hanya memuat satu grup. Orde lainnya
terbukti tak siklik seperti di atas: $4, 8, 9, 12$ oleh faktor
kuadrat, dan $6, 10, 12, 14$ oleh $2 \mid q - 1$. (Kebalikannya ---
$\gcd(n, \varphi(n)) = 1$ mengakibatkan $n$ siklik --- juga benar;
pertanyaan 7 membuktikan kasus tak trivial pertamanya, $n = pq$ dengan
$p \nmid q - 1$.)

\end{solution}
